\section{Teoría de Conjuntos}

\subsection{Conjuntos y Elementos}

	Un conjunto se puede entender de forma intuitiva como una colección bien definida de objetos, llamados \emph{elementos} o \emph{miembros} de tal conjunto.\footnote{Lee por ejemplo el artículo \href{https://www.britannica.com/science/set-theory/Axiomatic-set-theory}{Axiomatic set theory}. Para profundizar más, ve la conferencia \href{https://youtu.be/hLFit88zTYk}{¿Teoría de conjuntos? ¡Pero si es bien fácil!}, impartida en el Instituto de Matemáticas de la UNAM.}
	
	Usualmente, usaremos letras mayúsculas para denotar conjuntos, y minúsculas para los elementos. 

	Si un elemento $ x $ pertenece a un conjunto $ A $, escribiremos $ x\in A $. En caso contrario, $ x\not \in A $.
	El símbolo $\in$ significa \texttt{``es elemento de''}.  Por el contrario, $\notin$ significa \texttt{``no es elemento de''}.

	La pertenencia en un conjunto se denota de la siguiente manera:
	\begin{center}
		\emph{$a \in S$ denota que $a$ pertenece al conjunto $S.$} 
		
		\emph{$a,b \in S$ denota que tanto $a$ como $b$ pertenecen al conjunto $S.$}
	\end{center}
	
	Por ejemplo, el conjunto de los dígitos se puede describir como
	\[
		\set{\texttt{dígitos}}=\set{0,1,2,3,4,5,6,7,8,9}=\set{0,1,\dots,9}.
	\]

\subsection{Especificación de Conjuntos}

	Básicamente, existen dos maneras de especificar un conjunto en particular.  Por un lado, si es posible, enlistar todos los miembros (forma \emph{extensiva}).  Por otro lado, caracterizando los elementos en el conjunto (forma \emph{intensiva}).

	En cualquier caso, para declarar un conjunto se utilizan llaves:
	$$
	A=\set{\cdots}
	$$

	Por ejemplo, el conjunto 
	$$
	A=\set{1,3,5,6,9}
	$$
	también se puede especificar como
	$$
	A=\set{x\in \N \mid x<10, 2\nmid x }
	$$

	Un conjunto no depende del modo en que sus elementos se muestren.  Este permanece igual si sus elementos se repiten o se reacomodan.

	\begin{problema}
		\begin{align*}
			\set{x\in\R \mid x^{2}-3x+2=0} & =\set{1,2} \\
			&=\set{1,2,2,1}
		\end{align*}
	\end{problema}

\subsection{Símbolos especiales}

	Algunos conjuntos numéricos tienen una notación especial:
	\begin{itemize}
		\item $\N = \set{0,1,2,3,\dots}:$ números naturales (enteros no negativos); 
		\item $\Z = \set{0,\pm 1, \pm 2, \pm 3,\dots}:$ números enteros; 
		\item $\Z^{+} = \set{1,2,3,\dots}:$ enteros positivos;
		\item $\Q:$ números racionales; 
		\item $\R:$ números reales; 
		\item $\mathbb{C}:$ números complejos.
	\end{itemize}
	
	Observe que
	$$
	\N \subset \Z \subset \Q \subset \R \subset \mathbb{C},
	$$ pero en ningún caso los conjuntos son iguales.

\subsection{Subconjuntos}

	Supongamos que cada elemento en un conjunto $A$ es también elemento del conjunto $B,$ es decir,
	$$
	x\in A \onlyif x\in B.
	$$

	En ese caso, decimos que $A$ es subconjunto de $B.$  También podemos decir que $A$ está contenido en $B$ o que $B$ contiene a $A$.

	Esta relación se escribe como
	$$
	A \subset B 
	$$ o en ocasiones como
	$B\supset A.$  Por el contrario, si es necesario indicar que $A$ \emph{no} es  subconjunto de $B,$ escribimos $A \not\subset B.$

	Diremos que dos conjuntos son iguales si poseen los mismos elementos, es decir, 
	$$
	x \in A \iff x \in B.
	$$

	De manera equivalente 
	\begin{center}
		$A=B$ si y solo si $A \subset B$ y $B \subset A.$
	\end{center}

	\begin{problema}
		\label{lip:exmp:1.2}
		Determine la relación entre los siguientes conjuntos
		\begin{center}
			$$A=\set{1,3,4,7,8,9}, \, 
			B=\set{1,2,3,4,5}, \,
			C=\set{1,3}.$$
		\end{center}
	\end{problema}

	\begin{problema}
		Demuestre que 
		\begin{enumerate}
			\item $A\not\subset B$ si y solo si $\exists x\in A: x\notin B.$
			\item $A \subset A.$
			\item $A\subset B, B\subset C \onlyif  A\subset C.$
		\end{enumerate}
	\end{problema}

\subsection{Conjunto Universal y Conjunto Vacío}

	Todos los conjuntos bajo investigación en una aplicación de teoría de conjuntos se supone que pertenecen a un conjunto fijo más grande llamado \emph{conjunto universo} $\uset,$ a menos que se indique otro caso.

	Dado un conjunto universal $\uset$ y una propiedad $P,$ es posible que no existan elementos de $\uset$ con la propiedad $P.$ 

	Por ejemplo, el siguiente conjunto no tiene elementos
	$$
	S=\set{x\in \Z \mid x^{2}=3}.
	$$

	A tal conjunto sin elementos $\set{}$ se le conoce como conjunto vacío y se denota como $\emptyset.$

	\begin{observacion}
		\emph{Sólo existe un conjunto vacío}.  El conjunto vacío es subconjunto de cualquier otro conjunto.
	\end{observacion}

\subsection{Conjuntos disjuntos}

	Dos conjuntos $A$ y $B$ son \emph{disjuntos} si no tienen elementos en común, es decir, $A \cap B = \emptyset$.

	\begin{problema}
		Considere $$
		A=\set{1,2}, \; B=\set{4,5,6}, \; C=\set{5,6,7,8}.
		$$
		Determine qué pares de conjuntos son disjuntos. 
	\end{problema}

\subsection{Diagramas de Venn}

	Un diagrama de Venn es una representación gráfica de conjuntos en el que cada conjunto está representado por áreas encerradas en el plano.

	El conjunto universo $\uset$ es representado por el interior de un rectángulo, y cualquier otro conjunto está representado por discos que viven dentro del rectángulo.

	\begin{figure}
		\centering
		\includegraphics[width=11cm,keepaspectratio=true]{./md/venn01.png}
		\caption{Representaciones con Diagramas de Venn}
		\label{fig:0101}
	\end{figure}

\subsection{Unión e Intersección}

	La unión de dos conjuntos $A$ y $B$ es el conjunto de todos los elementos que pertenecen a $A$ o a $B,$  es decir
	$$
	A \cup B = \set{x \mid x\in A \vee x\in B}.
	$$

	La intersección de dos conjuntos $A$ y $B$ es el conjunto de todos los elementos que pertenecen a $A$ y a $B,$  es decir
	$$
	A \cap B = \set{x \mid x\in A \wed x\in B}.
	$$

	\begin{figure}
		\centering
		\includegraphics[width=10cm,keepaspectratio=true]{./md/venn_union_interseccion.png}
		\caption{Unión e Intersección}
		\label{fig:0103}
	\end{figure}

\subsection{Conjunto Potencia}

	Dado un conjunto $A,$ el \emph{conjunto potencia} de $A,$ denotado por $\mathcal{P}(A)$ o $2^{A},$ es el conjunto de \emph{todos} los subconjuntos de $A.$ Es decir,
	$$
	\mathcal{P}(A) = \set{X \mid X \subset A}.
	$$

	\begin{problema}
		Si $A = \set{a,b,c},$ entonces
		$$
		\mathcal{P}(A) = \set{\emptyset, \set{a}, \set{b}, \set{c}, \set{a,b}, \set{a,c}, \set{b,c}, \set{a,b,c}}.
		$$
		Note que $\card{\mathcal{P}(A)} = 8 = 2^{3}.$
	\end{problema}

	\begin{teorema}
		Si $A$ es un conjunto finito con $n$ elementos, entonces su conjunto potencia tiene $2^{n}$ elementos:
		$$
		\card{\mathcal{P}(A)} = 2^{\card{A}}.
		$$
	\end{teorema}

	\begin{proof}
		Cada elemento de $A$ puede estar o no estar en un subconjunto.  Como hay $\card{A}$ elementos y cada uno tiene 2 opciones (estar o no estar), el número total de subconjuntos es $2^{\card{A}}.$
	\end{proof}

\subsection{Teorema fundamental de subconjuntos}

	\begin{teorema}
		Sean $A$ y $B$ conjuntos.  Las siguientes proposiciones son equivalentes:
		\begin{enumerate}
			\item $A \subset B$
			\item $A \cap B = A$
			\item $A \cup B = B$
		\end{enumerate}
	\end{teorema}

	\begin{proof}
		Demostremos el ciclo $(1) \Rightarrow (2) \Rightarrow (3) \Rightarrow (1).$
		
		$(1) \Rightarrow (2):$ Supongamos $A \subset B.$  Claramente $A \cap B \subset A.$  Ahora, si $x \in A,$ entonces $x \in B$ por hipótesis, así que $x \in A \cap B.$  Por tanto $A \subset A \cap B,$ y concluimos $A \cap B = A.$
		
		$(2) \Rightarrow (3):$ Supongamos $A \cap B = A.$  Si $x \in B,$ entonces $x \in A \cup B.$  Si $x \in A,$ entonces $x \in A \cap B = A,$ así que $x \in B.$  En cualquier caso $x \in B,$ luego $A \cup B = B.$
		
		$(3) \Rightarrow (1):$ Supongamos $A \cup B = B.$  Si $x \in A,$ entonces $x \in A \cup B = B,$ así que $A \subset B.$
	\end{proof}
